\section{Input and Output Structure}

The user-facing input path is the Tkinter desktop workflow. A student can start from a blank model or shear-frame template, define nodes and members on the canvas, assign material and section properties, apply supports and loads, assign masses for modal analysis, validate the model, and then run Static or Modal analysis. The object tree and property panel support editing without writing XML or source code.

XML is the reproducibility backend. A desktop-created model can be saved/exported to XML and later loaded into the same \codeterm{ModelBuilder} and \codeterm{StructuralModel} path. The XML stores analysis-relevant geometry, materials, sections, supports, releases, loads, masses, unit settings, and load cases, while ordinary users work through the desktop UI.

Static outputs include the DOF map, \(K\), \(K_{ff}\), \(F\), \(F_f\), displacements, reactions, member-end forces, and N/V/M data. Modal outputs include \(K\), \(M\), reduced or condensed matrices, eigenvalues, frequencies, periods, mode shapes, modal masses, participation factors, effective modal masses, and mass participation ratios. Visible tables and plots can be exported where implemented. The installation manual covers environment setup, and the user manual gives a complete workflow example.
